\clearpage
\section{Fuentes informales relevantes}
En esta sección se registraron fuentes informales relevantes para complementar la búsqueda bibliográfica, identificar experiencias prácticas y reconocer tendencias relacionadas con la simulación de incendios forestales, la gestión de emergencias y el uso de herramientas geoespaciales. Estas fuentes se utilizaron como apoyo contextual y no como sustituto de los artículos científicos o documentos técnicos.
\begin{itemize}
    \item \textbf{Missoula Fire Sciences Laboratory / FRAMES --- Fire Research and Management Exchange System} (USDA Forest Service). Repositorio institucional que centraliza publicaciones, herramientas y documentación de los modelos de comportamiento del fuego (FARSITE, FlamMap, BehavePlus). Es la fuente primaria de los manuales técnicos y los datos de catalogación usados en esta bibliografía. \url{https://www.frames.gov/}
    \item \textbf{LANDFIRE} (USDA Forest Service y U.S. Department of the Interior). Portal de datos geoespaciales de vegetación, modelos de combustible, topografía y regímenes de fuego a 30 m de resolución. Es el tipo de insumo SIG que alimentaría la generación masiva de escenarios del proyecto. \url{https://landfire.gov/}
    \item \textbf{FDS and Smokeview --- Fire Dynamics Simulator} (National Institute of Standards and Technology, NIST). Sitio oficial del solucionador CFD sobre el que se construye WFDS, con guías de usuario, de verificación y de validación. \url{https://pages.nist.gov/fds-smv/}
    \item \textbf{Repositorio \texttt{firemodels/fds} en GitHub}. Código fuente, \emph{issues} y discusiones de FDS/WFDS y Smokeview. Es la vía práctica para conocer formatos de entrada/salida y limitaciones reales de la herramienta, información necesaria para construir el adaptador FARSITE\,$\rightarrow$\,WFDS. \url{https://github.com/firemodels/fds}
    \item \textbf{Wildland Fire Decision Support System (WFDSS NextGen)} (agencias federales de EE.\,UU.). Sistema operativo de apoyo a la decisión en incendios activos; sirve como referente funcional de lo que un comandante de incidente espera de una aplicación en campo. \url{https://wfdss.firenet.gov/}
    \item \textbf{International Association of Wildland Fire (IAWF)}. Asociación profesional internacional fundada en 1990 que edita el \emph{International Journal of Wildland Fire} y organiza congresos y \emph{summits} de seguridad en incendios forestales. \url{https://www.iawfonline.org/}
    \item \textbf{Centro de Estudos sobre Incêndios Florestais (CEIF/ADAI), Universidade de Coimbra}. Grupo organizador de la \emph{International Conference on Forest Fire Research} (ICFFR), celebrada cada cuatro años desde 1990; su décima edición tendrá lugar en Coimbra en 2026. Es el principal foro europeo sobre comportamiento del fuego e interfaz urbano-forestal. \url{https://www.uc.pt/adai/ceif/}
    \item \textbf{EFFIS --- European Forest Fire Information System} (Copernicus Emergency Management Service, Joint Research Centre de la Comisión Europea). Sistema operativo europeo de información sobre incendios: pronóstico de peligro, mapeo de áreas quemadas y estadísticas. Ejemplo de arquitectura de datos y servicios a escala continental. \url{https://forest-fire.emergency.copernicus.eu/}
    \item \textbf{Technosylva --- Wildfire Analyst}. Empresa que comercializa simulación de propagación en tiempo real e inteligencia de riesgo para agencias de bomberos y empresas de servicios públicos. Es el referente industrial más cercano a la aplicación propuesta y permite contrastar el aporte del proyecto frente al estado del arte comercial. \url{https://technosylva.com/}
    \item \textbf{Pyregence Consortium --- Open-Source Fire Science}. Consorcio de instituciones que desarrolla modelos de simulación y herramientas de pronóstico de incendios de código abierto (entre ellos ELMFIRE) para gestores de emergencias y comunidades. Referencia de cómo se empaqueta ciencia del fuego en herramientas operativas abiertas. \url{https://pyregence.org/}
\end{itemize}